%---------------------------------------------------------------------
\section{Finger Injury Management}
\label{sec:injury}

\begin{warningbox}{Evidence Caveat: No Climbing-Specific Pulley Rehabilitation RCTs}
No randomized controlled trials (RCTs) have been conducted on climbing pulley-tear rehabilitation protocols, return-to-climb criteria, or taping versus no-taping for clinical outcomes. Climbing-specific intervention evidence exists for hangboard training and is referenced where applicable, but pulley-specific rehabilitation protocols are classified \textbf{Bioplausible} (mechanistically extrapolated from non-climbing tendinopathy RCTs or tissue biology) or \textbf{Folklore} (practitioner/community heuristic). Evidence-backed claims are restricted to diagnostic imaging accuracy, biomechanics studies, the Schneeberger \& Schweizer (2016) conservative splint series, and the Sjöman et al.\ (2023) injury-risk survey. This limitation must be front-loaded in any clinical application of these protocols.
\end{warningbox}

%---------------------------------------------------------------------
\subsection*{Summary Table}

{\footnotesize
\setlength{\tabcolsep}{4pt}
\renewcommand{\arraystretch}{1.2}
\begin{longtable}{>{\raggedright\arraybackslash}p{2.0cm}
                  >{\raggedright\arraybackslash}p{1.8cm}
                  >{\raggedright\arraybackslash}p{1.8cm}
                  >{\raggedright\arraybackslash}p{2.6cm}
                  >{\raggedright\arraybackslash}p{2.8cm}
                  >{\raggedright\arraybackslash}p{1.2cm}
                  >{\raggedright\arraybackslash}p{1.8cm}}
\caption{Peer-reviewed studies cited in this section. S = surrogate; F = functional.}\label{tab:injury_studies}\\
\toprule
\textbf{Authors} & \textbf{Design / $n$} & \textbf{Injury type} & \textbf{Primary outcome} & \textbf{Effect size / key stat} & \textbf{S/F} & \textbf{Quality} \\
\midrule
\endfirsthead
\toprule
\textbf{Authors} & \textbf{Design / $n$} & \textbf{Injury type} & \textbf{Primary outcome} & \textbf{Effect size / key stat} & \textbf{S/F} & \textbf{Quality} \\
\midrule
\endhead
\bottomrule
\endfoot
\bottomrule
\endlastfoot

Sjöman et al.\ \cite{SGJ23}
& Cross-sectional survey; $n=434$ Finnish climbers
& Finger (SPIIF onset)
& Association of regular finger training (RFT) with sport-impairing injury by experience level
& $\geq$6\,yr: $\chi^2=4.57$, $p=0.03$ (protective); $<$6\,yr, $\geq$7a: $\chi^2=4.61$, $p=0.03$ (risk)
& \textbf{F} (injury incidence); observational; no causal inference
& Moderate \\

\addlinespace

Klauser et al.\ \cite{KFB+02}
& Diagnostic accuracy; $n=64$ climbers, 75 fingers
& Pulley injuries (A2/A3/A4)
& Dynamic US sensitivity/specificity vs.\ MRI reference
& Sensitivity 98\%, specificity 100\%; verification bias acknowledged
& \textbf{F} (diagnostic accuracy)
& Moderate \\

\addlinespace

Berrigan et al.\ \cite{BWC+21}
& SR imaging review
& A2 pulley diagnosis
& US vs.\ MRI/CT sensitivity/specificity
& Sensitivity/specificity 0.94--1.00
& \textbf{F} (diagnostic accuracy)
& Moderate; secondary synthesis \\

\addlinespace

Sch\"{o}ffl et al.\ \cite{SHWS03}
& Prospective registry cohort; $n=604$ climbers, 122 pulley injuries
& Pulley injuries Grade I--IV
& Return to prior climbing level after conservative vs.\ surgical management
& All conservative Gr.~I--III eventually returned; 7 had prolonged recovery due to persistent tenosynovitis; Grade IV surgical outcomes largely good; no controls
& \textbf{F} (return to sport); no ES; no RCT
& Low--moderate \\

\addlinespace

Schneeberger \& Schweizer \cite{SS16}
& Prospective case series; $n=47$ ruptures in 45 climbers
& Complete pulley ruptures (Gr.~II--IV)
& TBD reduction and return to prior climbing level with PPS
& A2 TBD: 4.4$\pm$1.0\,mm $\to$ 2.3$\pm$0.6\,mm (50\% reduction); 38/43 (88.4\%) regained prior level; mean return 8.8\,months
& \textbf{F} (TBD surrogate + return to sport)
& Moderate (prospective case series; no controls) \\

\addlinespace

Mohn et al.\ \cite{MSM+22}
& Retrospective descriptive; $n=65$ climbers
& Flexor tenosynovitis
& Return to climbing and symptom duration
& All returned to climbing; ${\approx}$75\% regained or exceeded prior level; mean symptom duration 30.5\,wk; outcomes independent of specific therapy content
& \textbf{F} (return to sport)
& Low--moderate (retrospective; no controls) \\

\addlinespace

Lutter et al.\ \cite{LSS+18}
& Consecutive case series; $n=60$ (57 climbers)
& Lumbrical muscle tear
& Return to climbing by grade; imaging classification
& All recovered with functional therapy; Grade III healing significantly longer; return 6$+$\,wk (Gr.~I--II), 3--6\,months (Gr.~III)
& \textbf{F} (return to sport); imaging by US/MRI
& Moderate (largest dedicated series; no controls) \\

\addlinespace

Sch\"{o}ffl et al.\ \cite{SLL+25}
& Prospective case series; $n=69$ capsulitis cases
& PIP joint capsulitis/synovitis
& Return to prior climbing level at 12\,months
& 72\% returned to prior level; 28\% showed decreased performance (mainly from extended rest, not residual symptoms); staged algorithm effective
& \textbf{F} (return to sport)
& Moderate (prospective case series; no controls) \\

\addlinespace

Schweizer \cite{Sch00}
& Biomechanical in-vivo; $n=16$ fingers
& A2 pulley loading / taping
& Bowstringing distance and force absorption under taping variants
& Proximal A2: bowstringing $\downarrow$2.8\%, force absorbed ${\approx}11\%$; distal: $\downarrow$22\%, ${\approx}12\%$
& \textbf{S} (biomechanical proxy; no clinical outcome)
& Moderate (biomechanics) \\

\addlinespace

Alfredson et al.\ \cite{APJL98}
& Prospective cohort $+$ historical controls; $n=15+15$
& Chronic Achilles tendinosis (analogue)
& Return to running at 12\,wk: eccentric heel-drop vs.\ continued activity
& 100\% return (eccentric) vs.\ 0\% (control midpoint); no formal ES
& \textbf{F} (return to sport); transfer to pulleys bioplausible only
& High (Achilles); ESS 3 pulley transfer \\

\addlinespace

Cook \& Purdam \cite{CP08}
& Conceptual review / mechanistic model
& Tendinopathy (general)
& Reactive$\to$dysrepair$\to$degenerative model; load-management framework
& No ES; conceptual only
& \textbf{S} (mechanistic model)
& Moderate; well-cited \\

\addlinespace

Rio et al.\ \cite{RKP+15}
& RCT crossover; $n=6$/condition (patellar tendinopathy)
& Patellar tendinopathy (analogue)
& Immediate pain (VAS) and MVIC: isometric vs.\ isotonic
& Pain: 7.0$\pm$2.04 $\to$ 0.17$\pm$0.41 (isometric); MVIC $\uparrow$18.7$\pm$7.8\%; no CI reported
& F (VAS pain); S (MVIC); patellar only
& Moderate; $n=6$; single session \\

\addlinespace

Shaw et al.\ \cite{SLBR+17}
& RCT crossover; $n=8$
& Connective tissue (ligament model)
& P1NP AUC; engineered-ligament mechanics post-supplement
& P1NP ${\approx}2{\times}$ placebo ($p{<}0.05$); ligament stiffness $\uparrow$; no CI reported
& \textbf{S} (P1NP; in-vitro ligament)
& Moderate; $n=8$ \\

\addlinespace

Iruretagoiena et al.\ \cite{IUDORL+20}
& SR; 7 studies
& A2/A4 pulley diagnosis
& US tendon--bone distance cutoff thresholds
& No consensus; A2 rupture cutoffs 1.9--5.1\,mm across studies
& \textbf{F} (diagnostic accuracy); threshold heterogeneity limits utility
& Moderate diagnostic SR \\

\end{longtable}
}% end footnotesize

%---------------------------------------------------------------------
\subsection{Injury-Risk Stratification Framework}
\label{subsec:riskstrata}

Risk tiers used throughout this section. Incidence estimates derive from available survey and case series data; where climbing-specific data are absent, estimates are flagged as extrapolated from adjacent sports.

\renewcommand{\arraystretch}{1.3}
\begin{center}
\small
\begin{tabular}{>{\raggedright}p{2.4cm}
                >{\raggedright}p{3.0cm}
                >{\raggedright}p{4.0cm}
                >{\raggedright\arraybackslash}p{5.0cm}}
\toprule
\textbf{Stratum} & \textbf{Estimated incidence} & \textbf{Definition} & \textbf{Management implication} \\
\midrule
\textbf{Low} & $<$5\,\%/year & Managed with standard warm-up and load monitoring; no elevated structural risk & Routine load progression; standard 48\,h recovery; no additional screening \\
\addlinespace
\textbf{Moderate} & 5--20\,\%/year & Context-dependent risk; structural risk present but mitigable with specific precautions & Explicit monitoring of finger pain VAS; HRV supplementary tracking; grip-position modification if VAS $\geq$2/10 \\
\addlinespace
\textbf{High} & $>$20\,\%/year & Elevated structural risk; explicit mitigation required & Mandatory load reduction, grip-position restriction, structured rehabilitation; medical evaluation if VAS $\geq$3/10 at rest \\
\addlinespace
\textbf{Very High} & Near-certain injury without intervention & Contraindicated without medical supervision & Full climbing rest pending medical evaluation; return protocol requires clinical clearance \\
\bottomrule
\end{tabular}
\end{center}

\noindent\textbf{Note on incidence data.} No prospective climbing study has measured annual pulley injury incidence by stratum. The rate estimates above are informed by Sch\"{o}ffl et al.\ (2003, \cite{SHWS03}) and Sjöman et al.\ (2023) survey data and should be treated as indicative. Injury incidence and return-to-sport risk are separate quantities: a Low-incidence injury can carry High return-to-sport risk if mismanaged. These strata are stated explicitly at each recommendation below where applicable.

%---------------------------------------------------------------------
\subsection{A2 Pulley Injuries}
\label{subsec:a2pulley}

\subsubsection{Anatomy and Mechanism}

The flexor tendon sheath of digits II--V contains five annular (A1--A5) and three cruciate (C1--C3) pulleys. The A2 pulley spans the proximal half of the proximal phalanx and is the largest, least deformable annular pulley; it prevents bowstringing of the flexor digitorum profundus (FDP) and flexor digitorum superficialis (FDS) by maintaining the tendons' mechanical line of action. In published case series, the ring (digit IV) and middle (digit III) fingers account for $\approx$98\% of pulley injuries in climbers, consistent with their role as the primary force-bearing digits in crimp configurations.

\begin{bioplausbox}{Full-Crimp Loading Hierarchy and Eccentric Failure Mechanism · \bioplausible · ESS 2}
\textbf{Grip-position loading hierarchy (Bioplausible):} Full crimp (PIP $\approx$90$^\circ$ flexion, DIP hyperextension or 0--20$^\circ$) $>$ half crimp $>$ open hand, due to acute tendon approach angle against the pulley increasing contact pressure. Schweizer (2001) estimated full-crimp A2 forces at 35--40\,N/kg body mass in cadaveric/model studies, half-crimp at $\approx$30--35\,N/kg, and open-hand at $\approx$20--25\,N/kg. No 95\,\% CI was reported. The Vigouroux et al.\ (2006, \cite{VQLVM06}) biomechanical model reports A2 forces $\approx$36$\times$ higher in crimp versus slope grip, and A4 forces $\approx$4$\times$ higher --- a substantially larger differential than the Schweizer absolute-force figures alone convey.

Sch\"{o}ffl I.\ et al.\ (2009, \cite{SOJ+09}) cadaveric study of concentric vs.\ eccentric loading found \textit{eccentric} pulley stress (sudden finger extension under load, as in a foot slip or a hold crimped open by surprise) to be the dominant failure mechanism, not static contact pressure alone. These data support recommending open-hand grip during volume accumulation and caution during eccentric-loaded holds; however, no prospective injury-prevention RCT has confirmed that grip-position prescription reduces pulley injury incidence in climbers. \textbf{[Bioplausible: mechanistic from Schweizer 2000, Vigouroux 2006 \cite{VQLVM06}, and Sch\"{o}ffl 2009 \cite{SOJ+09}; not confirmed in prospective trial.]}
\end{bioplausbox}

\subsubsection{Grading}

The Sch\"{o}ffl (2003) grading system \cite{SHWS03} remains the dominant clinical and research classification:

\begin{itemize}[leftmargin=1.5em]
  \item \textbf{Grade I}: Pulley strain without significant tendon--bone dehiscence; TBD increase $<$1--2\,mm on dynamic US. Tenderness over pulley; pain with resisted PIP flexion.
  \item \textbf{Grade II}: Complete rupture of A4 pulley \textit{or} partial rupture of A2 or A3 pulley. Mild or absent bowstringing; palpable swelling; pain with active loading.
  \item \textbf{Grade III}: Complete rupture of A2 \textit{or} A3 pulley. Clear bowstringing (TBD $>$2--3\,mm under load on US); may have bruising and palpable defect.
  \item \textbf{Grade IVa}: Double-pulley rupture (A2/A3 or A3/A4). Substantial bowstringing; significant functional deficit (Simon et al.\ 2022, \cite{SLTS22}).
  \item \textbf{Grade IVb}: Triple-pulley rupture (A2/A3/A4) \textit{or} single/double rupture with complicating pathology: Flap Irritation Phenomenon (FLIP), progressive contracture, tenosynovitis, collateral ligament injury, or lumbrical involvement (\cite{SLTS22}).
\end{itemize}

\begin{bioplausbox}{Schöffl Grading System: Clinical Utility Without RCT Validation · \bioplausible · ESS 2}
\textbf{Grading validity (Bioplausible):} The Sch\"{o}ffl system derives from case series, not a validated prospective grading-outcome study. No inter-rater reliability data from RCTs exists. Grade IV was refined by Simon et al.\ (2022, \cite{SLTS22}) into IVa (double-pulley) and IVb (triple-pulley or with complications), which has direct surgical-threshold implications. The system is clinically useful for mapping structural severity to management pathway but should not be treated as a validated prognostic instrument. \textbf{[Bioplausible: supported by case series \cite{SHWS03}; Grade IV refinement from Simon et al.\ 2022 \cite{SLTS22}; no RCT validation.]}
\end{bioplausbox}

\subsubsection{Diagnosis}

\paragraph{Physical examination.}
Pain reproduced by resisted PIP flexion and direct palpation over the proximal volar aspect of the proximal phalanx is the primary clinical indicator. Bowstringing may be palpable or visible in Grade III--IV when the digit is actively flexed against resistance. Sensitivity and specificity of clinical examination alone are not established in controlled diagnostic accuracy studies.

\begin{bioplausbox}{Clinical Examination Reliability for Pulley Injury: Unvalidated · \bioplausible · ESS 2}
\textbf{Clinical exam reliability (Bioplausible):} Focal tenderness and pain with resisted PIP flexion are widely used but unvalidated against imaging or surgical reference standards in controlled studies. \textbf{[Bioplausible: clinical reasoning; no ROC/diagnostic accuracy study.]}
\end{bioplausbox}

\paragraph{Imaging.}\leavevmode

\begin{evidencebox}{Ultrasound as Primary Diagnostic Modality for Pulley Injury · \evidencebacked · ESS 6}
\textbf{Ultrasound as primary modality (Evidence-backed):} Dynamic high-resolution US ($\geq$15\,MHz linear probe, assessed in passive extension and full resisted flexion) is the recommended first-line modality. Klauser et al.\ (2002) reported US sensitivity 98\%, specificity 100\% vs.\ MRI reference standard in 64 climbers with 75 symptomatic fingers (\cite{KFB+02}; verification bias acknowledged). Berrigan et al.\ (2021) systematic review confirms US sensitivity/specificity 0.94--1.00 across included studies (\cite{BWC+21}). Dynamic assessment comparing TBD at 20$^\circ$ passive extension vs.\ full active flexion is essential; static views underestimate Grade II bowstringing and miss partial tears.

Complete A2 rupture TBD diagnostic cutoffs are inconsistent across studies, ranging 1.9--5.1\,mm (Iruretagoiena et al.\ 2020, \cite{IUDORL+20}). Miro et al.\ (2021, \cite{MVSS21}) found that experienced high-level climbers have physiologically elevated TBD (0.3--0.4\,mm longer at A2/A4) versus non-climbers \textit{without reported injury}, indicating that TBD thresholds derived from non-climbing populations may over-diagnose injury in adapted athletes. MRI sensitivity for complete pulley tears ranges 0.79--1.00 across studies in varying finger positions; MRI is superior for multi-structure injuries (capsulitis, collateral ligament, volar plate). \textbf{Bottomline: use high-resolution dynamic US first; adjust TBD thresholds for climber adaptation; reserve MRI for equivocal presentations or suspected multi-structure Grade IVa/IVb injury.}
\end{evidencebox}

%---------------------------------------------------------------------
\subsubsection{Conservative Management --- Grade I--II and Most Grade III}

\begin{warningbox}[title=No Climbing Rehabilitation RCTs]
All immobilisation durations, return-to-loading timelines, and hangboard progression thresholds below are derived from case series (Sch\"{o}ffl et al.\ 2003, \cite{SHWS03}) or extrapolated from non-climbing tendinopathy RCTs. None has been validated in a controlled climbing study.
\end{warningbox}

\paragraph{Immobilisation.}\leavevmode

\begin{bioplausbox}{Short-Term Immobilisation: Grade-Based Duration (Schöffl 2003) · \bioplausible · ESS 2}
\textbf{Short-term immobilisation (Bioplausible):} Typical recommendations from Sch\"{o}ffl et al.\ (2003, \cite{SHWS03}) case series:
\begin{itemize}[leftmargin=1.5em,noitemsep]
  \item Grade I: 0\,d immobilisation; functional therapy 2--4\,wk; tape protection.
  \item Grade II: 10\,d immobilisation with pulley-protection splint (PPS); then 2--4\,wk functional therapy.
  \item Grade III: 14\,d immobilisation with PPS; PPS worn continuously for 6--8\,wk total; functional therapy alongside.
\end{itemize}
Rationale: early controlled mechanical stimulation promotes aligned collagen deposition; prolonged immobilisation ($>$1\,wk) is inconsistent with tendon and ligament healing biology. However, no climbing RCT has compared immobilisation durations. \textbf{[Bioplausible: extrapolated from tendon biology; case series source.]}
\end{bioplausbox}

\begin{evidencebox}{Pulley-Protection Splint (PPS): Primary Conservative Device for Grade II--III · \evidencebacked · ESS 5}
\textbf{Pulley-protection splint (Evidence-backed for conservative outcomes):} The PPS (thermoplastic circumferential ring worn full-time) physically approximates the tendon to the bone continuously, reducing TBD passively at rest and during activity. Schneeberger \& Schweizer (2016, \cite{SS16}) prospective series ($n=47$ complete ruptures in 45 climbers):
\begin{itemize}[leftmargin=1.5em,noitemsep]
  \item PPS worn continuously for 2\,months; active functional therapy alongside from early weeks.
  \item A2 TBD reduced from 4.4$\pm$1.0 to 2.3$\pm$0.6\,mm (50\% reduction); A4 TBD from 2.9$\pm$0.7 to 2.1$\pm$0.5\,mm.
  \item 38/43 climbers (88.4\%) regained prior climbing level at mean 8.8\,months.
\end{itemize}
The Sch\"{o}ffl (2020, \cite{SSF+20}) MLTJ review states: ``The pulley support ring therapy must be started within a few days after trauma and performed strictly for 6--8 weeks.'' PPS and H-tape serve complementary roles: PPS provides continuous passive bowstringing reduction during healing; H-tape provides additional support during sport. \textbf{[Evidence-backed for TBD reduction and functional return: prospective case series \cite{SS16}; no RCT comparison vs.\ H-tape or immobilisation only.]}
\end{evidencebox}

\paragraph{Return-to-loading timeline.}\leavevmode

\begin{bioplausbox}{Return-to-Climbing Timelines After Pulley Injury · \bioplausible · ESS 2--5}
\textbf{Return-to-climbing timelines:} Sch\"{o}ffl et al.\ (2003, \cite{SHWS03}) case series provides the foundational reference:
\begin{itemize}[leftmargin=1.5em,noitemsep]
  \item Grade I: progressive reloading at 1--2\,wk; return to climbing 4--6\,wk.
  \item Grade II: progressive reloading at 3--4\,wk; return to climbing 6--10\,wk.
  \item Grade III (conservative with PPS): progressive reloading at 6--8\,wk; Schneeberger \& Schweizer (2016, \cite{SS16}) prospective PPS series reports mean 8.8\,months to prior climbing level --- substantially longer than the Sch\"{o}ffl (2003) ``3--5\,months'' figure, which derived from a subset without continuous PPS. Grade III patients should be counselled on 6--10\,months as a realistic expectation.
\end{itemize}
\textbf{[Grade I--II: Bioplausible/Folklore from Sch\"{o}ffl 2003 \cite{SHWS03}. Grade III: ESS 5 prospective case series \cite{SS16}; mean 8.8\,months.]}
\end{bioplausbox}

\paragraph{H-taping and ring splinting.}\leavevmode

\begin{bioplausbox}{H-Taping: Biomechanical Evidence and Limitations · \bioplausible · ESS 2}
\textbf{Biomechanical basis and limitations (Bioplausible):} Schweizer (2000, \cite{Sch00}) in-vivo study ($n=16$ fingers) found tape placed distally over the proximal phalanx (the more effective position) reduced bowstringing by 22\% and absorbed $\approx$12\% of bowstringing force; more proximal placement over the A2 pulley gave only 2.8\% reduction. Sch\"{o}ffl et al.\ (2007, \cite{SES+07}) controlled biomechanical study ($n=8$ climbers) found H-tape reduced TBD by 16\% and increased crimp MVC by $+$13\% (no CI) vs.\ no tape; no effect in open-hand hang.

A systematic review by Larsson et al.\ (2022, \cite{LNB22}) reported low-to-moderate certainty for bowstringing reduction (15--22\%) in \textit{injured} fingers only; no significant effect in uninjured fingers; zero RCTs on re-injury. Salas et al.\ (2022, \cite{SMT+23}) cadaveric biomechanical study (56 paired fresh-frozen specimens) found H-taping did \textit{not} significantly increase failure force or reduce bowstringing in partial A2 tears, and provided no protection against A2 rupture in intact specimens. The cadaveric findings directly challenge H-tape's biomechanical protective rationale; clinical utility may rest primarily on pain modulation, proprioception, and confidence. \textbf{[Bioplausible: positive small-sample data from \cite{Sch00,SES+07}; contradicted by cadaveric evidence \cite{SMT+23}; no RCT for clinical outcomes.]}
\end{bioplausbox}

\begin{folklorebox}{H-Taping in Climbing Rehabilitation: Universal Practice · \folkloreverified · ESS 1}
\textbf{H-taping clinical practice (Folklore):} H-taping around the proximal phalanx (figure-8 configuration creating dorsal bridge) is near-universal in climbing rehabilitation practice. Community consensus, Hooper's Beta, Schöffl-era clinical guidelines. No RCT or controlled case series has measured re-injury rates, healing time, or functional outcomes with vs.\ without taping. \textbf{[Folklore: climbing community / climbing physio consensus.]}
\end{folklorebox}

\paragraph{Load-managed rehabilitation --- tendinopathy extrapolation.}\leavevmode

\begin{evidencebox}{Tendon Loading Paradigms: Eccentric and Isometric Evidence (Achilles/Patellar) · \evidencebacked · ESS 3}
\textbf{Tendon loading paradigms (Evidence-backed for analogous tissues):}
\begin{itemize}[leftmargin=1.5em]
  \item Alfredson et al.\ (1998, \cite{APJL98}): heavy-load eccentric heel-drop (3$\times$15 twice daily, 12\,wk) achieved 100\% return to running in $n=15$ chronic Achilles tendinosis patients vs.\ 0\% for continued activity alone. A 2023 meta-analysis of 5 Achilles RCTs reported eccentric exercise pain reduction MD $-$1.21 (95\,\% CI $-$2.72 to $-$0.30) vs.\ other exercise.
  \item Cook \& Purdam (2009, \cite{CP08}): reactive--dysrepair--degenerative continuum model; reactive stage warrants load reduction (not rest) and isometric loading; degenerative stage benefits from progressive isotonic loading.
  \item Rio et al.\ (2015, \cite{RKP+15}): isometric leg extension produced immediate pain reduction of 6.8/10 (from 7.0$\pm$2.04 to 0.17$\pm$0.41, sustained 45\,min) and MVIC $\uparrow$18.7$\pm$7.8\% vs.\ isotonic in patellar tendinopathy ($n=6$).
\end{itemize}
\textbf{[Evidence-backed for Achilles and patellar tendinopathy. Extrapolation to A2 pulley is Bioplausible only --- see below.]}
\end{evidencebox}

\begin{bioplausbox}{Progressive Loading Applied to A2 Pulley: Extrapolation from Tendinopathy · \bioplausible · ESS 2}
\textbf{Pulley-specific loading (Bioplausible):} The A2 pulley is a fibrocartilaginous ligamentous condensation of the flexor sheath --- primarily type\,I collagen with fibrocartilaginous zones under compressive load --- not a tendon. Cellular overlap with tendons exists (fibroblast mechanosensing, integrin-mediated collagen remodelling), but load-response properties are not identical. Application of progressive loading protocols (isometrics $\to$ isotonic $\to$ functional) to pulley rehabilitation is bioplausible but constitutes an extrapolation from a different tissue type. No peer-reviewed study has measured A2 pulley load during hangboard exercises with validated force-sensing systems and correlated measurements with tissue healing outcomes. \textbf{[Bioplausible: mechanistic framework from Cook \& Purdam 2009 and Alfredson 1998; Folklore: specific thresholds and protocols.]}
\end{bioplausbox}

\paragraph{Load-managed hangboard rehabilitation during recovery.}\leavevmode

\begin{folklorebox}{Tyler Nelson Progressive Loading Protocols for Pulley Rehabilitation · \folkloreverified · ESS 1}
\textbf{Tyler Nelson / Camp 4 Human Performance protocols (Folklore):} Tyler Nelson has disseminated progressive loading protocols for pulley rehabilitation through clinical practice and online/professional channels. Proposed structure: (a) isometric hangs beginning at 20--40\% perceived maximum in the early reactive phase, pain $\leq$2--3/10; (b) progress load 5--10\% per week guided by symptom response; (c) grip-position progression --- open-hand $\to$ half crimp, full crimp avoided until $\approx$week 12. These thresholds are not empirically derived from pulley-specific loading studies. \textbf{[Folklore: Tyler Nelson, C4HP, practitioner/coaching channels. Mechanistic framework Bioplausible via Cook \& Purdam 2009; specific thresholds unvalidated.]}
\end{folklorebox}

%---------------------------------------------------------------------
\subsubsection{Surgical Management --- Grade III--IV}

\begin{bioplausbox}{Surgical Indications for Pulley Injury · \bioplausible · ESS 2--3}
\textbf{Surgical indications (Bioplausible / case series):} Current consensus (Simon et al.\ 2022, \cite{SLTS22}; Schneeberger \& Schweizer 2016, \cite{SS16}) classifies Grade I--III as conservative first-line. Surgery is indicated for:
\begin{itemize}[leftmargin=1.5em,noitemsep]
  \item Grade IVa or IVb with persistent bowstringing or functional deficit after adequate conservative management.
  \item Flap Irritation Phenomenon (FLIP): pulley stump folding under the intact sheath, producing constant mechanical friction; US shows oscillating hyperechoic structure; MRI confirms. Resection $\pm$ repair resolves (Sch\"{o}ffl et al.\ 2011, \cite{SBS11}).
  \item Grade III with documented failure of $\geq$6--9\,months of compliant PPS + functional therapy.
\end{itemize}
Techniques: graft-based loop reconstruction (palmaris longus, extensor retinaculum) remains the standard for Grade IVa/b multi-pulley loss. Direct suture repair without graft (Balandier et al.\ 2026, \cite{BBM+26}) was reported in 8 elite climbers (6 isolated, 2 multiple ruptures): mean return to pre-injury level 5.8\,months (range 3--8), US-confirmed pulley continuity at 3\,years --- a substantially shorter timeline than graft reconstruction but limited to small-$n$ series.

Arora et al.\ (2007) retrospective series ($n=23$) reported excellent/good Buck-Gramcko outcomes in 21/23.\singlesource Bouyer et al.\ (2016) retrospective series ($n=38$, mean FU 85\,months)\singlesource found post-operative US E-space $<$2\,mm predicted return in 70\% vs.\ 25\%. Return-to-climbing: 3--6\,months (Grade III failed conservative, direct suture); 6--12\,months (Grade IVa/b graft).

\textbf{[Bioplausible / Low-quality evidence: all data from case series; no controls; heterogeneous techniques. Grade III should not receive first-line surgery \textemdash{} Schneeberger \& Schweizer \cite{SS16} showed 88.4\% return with conservative PPS even for complete ruptures.]}
\end{bioplausbox}

%---------------------------------------------------------------------
\subsection{Flexor Tendon Strains}
\label{subsec:flexortendon}

\begin{bioplausbox}{Flexor Tendon Strain: Anatomy and Mechanism · \bioplausible · ESS 2}
\textbf{Anatomy and mechanism (Bioplausible):} FDP inserts distally to the distal phalanx; FDS splits at the proximal phalanx to insert bilaterally onto the middle phalanx base. Both run through the same fibrous sheath as the A2 pulley. Flexor tendon strains (partial or microtear injuries to the tendon proper) typically occur at the musculotendinous junction or within the tendon substance under repeated high-load crimp grip loading or sudden eccentric overload. They co-occur with Grade IV pulley injuries (definitional) or may present independently. No climbing-specific epidemiological data distinguish flexor tendon strain prevalence from pulley injury prevalence in isolation. \textbf{[Bioplausible: anatomical reasoning and clinical report.]}
\end{bioplausbox}

\paragraph{Grading and differentiation from pulley injury.}\leavevmode

\begin{bioplausbox}{Flexor Tendon Strain Grading and Differentiation from Pulley Injury · \bioplausible · ESS 2}
\textbf{Grading (Bioplausible / clinical consensus):}
\begin{itemize}[leftmargin=1.5em,noitemsep]
  \item Grade I: minor fibre disruption; no architectural loss on US; tenderness along tendon course; pain with resisted flexion.
  \item Grade II: partial tear; hypoechoic defect on US; weakened active flexion; focal swelling.
  \item Grade III: complete/near-complete disruption; functional deficit; requires urgent surgical referral.
\end{itemize}
Distinguishing from pulley injury: pulley injury presents with focal tenderness over the proximal phalanx with bowstringing on dynamic US; flexor tendon strain presents with pain along the tendon course proximal or distal to the pulley, internal tendon changes on US, but no bowstringing. Precise attribution requires high-quality dynamic US. \textbf{[Bioplausible: no climbing-specific grading validation study.]}
\end{bioplausbox}

\paragraph{Management.}\leavevmode

\begin{bioplausbox}{Grade I--II Flexor Tendon Strain: Conservative Management · \bioplausible · ESS 2}
\textbf{Grade I--II management (Bioplausible):} Principles mirror A2 pulley Grade I--II management: (a) acute load reduction 5--14\,d; (b) early progressive loading from week 2--3, pain-guided; (c) progressive return to full loading over 6--12\,wk. Mechanistic basis is identical to the Cook--Purdam continuum framework. No climbing-specific RCT or case series specifically distinguishes flexor tendon strain outcomes from pulley injury outcomes. \textbf{[Bioplausible: same extrapolation from tendinopathy literature.]}
\end{bioplausbox}

Grade III FDP avulsion (``jersey finger'' equivalent): requires urgent surgical referral and formal hand surgery protocol --- outside scope of conservative climbing rehabilitation.

%---------------------------------------------------------------------
\subsection{Flexor Tenosynovitis}
\label{subsec:tenosynovitis}

Flexor tenosynovitis --- inflammation of the tendon sheath surrounding the FDP/FDS tendons --- is the second most common climbing finger injury (7\% of all injuries in Sch\"{o}ffl 2003 \cite{SHWS03} and Lutter 2020 \cite{SSF+20} cohorts). It is the injury most frequently confused with pulley injury; mismanagement (treating as a pulley strain) prolongs recovery.

\begin{evidencebox}{Flexor Tenosynovitis: Diagnosis and Conservative Outcomes · \evidencebacked · ESS 5}
\textbf{Diagnosis:} US shows the \textit{halo phenomenon} --- peritendinous fluid in the tendon sheath around the flexor tendon. No bowstringing. Pain is diffuse along the tendon course (proximal to the pulley or radiating into the palm) rather than the focal proximal-phalanx tenderness of a pulley injury. Active glide against resistance reproduces pain.

\textbf{Natural history and conservative management} (Mohn et al.\ 2022, \cite{MSM+22}; retrospective $n=65$ climbers with confirmed tenosynovitis):
\begin{itemize}[leftmargin=1.5em,noitemsep]
  \item All patients returned to climbing; ${\approx}$75\% regained or exceeded prior climbing level.
  \item Mean symptom duration 30.5\,wk; course was undulating (flare-ups common).
  \item Load reduction was the dominant treatment (91\% of patients); no specific therapy content significantly predicted outcome in regression analysis, suggesting a favourable natural history.
  \item Predictors of longer recovery: high pain during daily activities at presentation; prior climbing grade $\geq$7b.
\end{itemize}

\textbf{Corticosteroid injection for chronic tenosynovitis} ($>$6\,wk; Sch\"{o}ffl et al.\ 2019, \cite{SSL19}; uncontrolled case series $n=42$): two injections 7--10\,days apart; 73.8\% pain-free after second injection; mean 84.2\% pain reduction; no complications. No surgery required in any case.

\textbf{[Evidence-backed for natural history and corticosteroid utility: case series \cite{MSM+22,SSL19}; no RCT; no control group.]}
\end{evidencebox}

%---------------------------------------------------------------------
\subsection{PIP Joint Capsulitis and Synovitis}
\label{subsec:capsulitis}

PIP joint capsulitis/synovitis is the third most common climbing finger injury (6--7\% of all injuries; $\approx$18.8\% of finger injuries per Sch\"{o}ffl 2020 \cite{SSF+20}). It is frequently misdiagnosed as a pulley injury because both present with PIP-region pain.

\begin{evidencebox}{PIP Capsulitis/Synovitis: Mechanism, Diagnosis, and Treatment Algorithm · \evidencebacked · ESS 5}
\textbf{Mechanism:} Repeated crimp loading generates high peak pressure in the PIP joint, driving inflammatory mediators, synovial effusion, and progressive capsular fibrosis. Untreated chronic capsulitis leads to early osteoarthritis; long-term climbers commonly show PIP osteophytes on imaging.

\textbf{Differential from pulley injury:} Diffuse PIP joint swelling and morning stiffness without bowstringing; pain on active flexion/extension throughout the arc (not just at load); absent focal tenderness over the proximal phalanx. US shows intra-articular effusion $\pm$ periarticular hyperemia.

\textbf{Staged treatment algorithm} (Sch\"{o}ffl et al.\ 2025, \cite{SLL+25}; prospective case series $n=69$; Tanguay et al.\ 2023, \cite{Vag23}):
\begin{itemize}[leftmargin=1.5em,noitemsep]
  \item \textbf{Stage 1}: Load reduction, anti-inflammatory (topical/oral short-course), contrast baths, joint mobilisation; avoid full crimp until resolved.
  \item \textbf{Stage 2} (failure of Stage 1 at 4--6\,wk): Corticosteroid or hyaluronic acid injection.
  \item \textbf{Stage 3} (chronic, refractory): Radiosynoviorthesis.
\end{itemize}
Outcome: 72\% returned to prior climbing level; 28\% showed decreased performance, predominantly attributed to prolonged forced rest rather than residual symptoms. All climbers resumed some climbing within 12\,months.

\textbf{[Evidence-backed: prospective case series \cite{SLL+25}; no RCT; no control condition.]}
\end{evidencebox}

%---------------------------------------------------------------------
\subsection{Lumbrical Muscle Injuries}
\label{subsec:lumbrical}

Lumbrical muscle tears are highly climbing-specific and are absent from most general sports medicine references. They are caused by the \textit{quadriga effect}: the 3rd and 4th lumbrical muscles have a unique bipennate dual-FDP origin; when the ring or middle finger is loaded in a pocket (flexed) while adjacent fingers are extended, a shearing force tears the lumbrical at its FDP attachment. Clinical presentation: pain in the \textit{palm} at or proximal to the MCP joint, radial side of the ring finger --- not over the proximal phalanx as in pulley injury.

\begin{evidencebox}{Lumbrical Tear: Diagnosis, Classification, and Outcomes · \evidencebacked · ESS 5}
\textbf{Diagnostic test (Lumbrical Stress Test):} Pain provoked by passive extension of the ring finger while the patient attempts to close their fist; pain disappears when the middle finger is simultaneously extended (neutralising the quadriga effect). This test is specific to lumbrical involvement.

\textbf{Classification} (Lutter et al.\ 2018, \cite{LSS+18}; $n=60$ consecutive patients, 57 climbers):
\begin{itemize}[leftmargin=1.5em,noitemsep]
  \item Grade I: microtrauma, US-negative; Grade II: muscle-fibre disruption visible on US; Grade III: musculotendinous disruption (US or MRI required).
\end{itemize}
\textbf{Outcomes:} All 60 patients recovered with adapted functional therapy. Grade III healing significantly longer than Grade I/II. Return to climbing: $\geq$6\,wk (Gr.~I--II); 3--6\,months (Gr.~III).

\textbf{Management:} Avoid single-finger pocket holds and differential-loading configurations during recovery. Progressive open-hand loading on matched holds; gradual reintroduction of pockets guided by pain response. Schweizer (2003, \cite{Sch03}) first described the clinical entity in three case reports.

\textbf{[Evidence-backed for natural history and classification: largest dedicated series \cite{LSS+18} ($n=60$); no RCT; no control.]}
\end{evidencebox}

%---------------------------------------------------------------------
\subsection{Adolescent Physeal Injuries}
\label{subsec:adolescent}

\begin{warningbox}{Adolescent Climbers: Pulley Framework Does Not Apply}
In climbers under $\approx$16--17 years (male) or $\approx$14--15 years (female), the PIP joint growth plate physis is the mechanical weak point, not the pulley. Crimp-grip loading preferentially stresses the open physis, producing \textit{primary periphyseal stress injuries} (PPSI). Presentation mimics a pulley strain (PIP swelling, pain with loading) but X-ray $\pm$ MRI is required. Sch\"{o}ffl et al.\ (2022, \cite{SSF+21}) prospective multicenter algorithm study ($n=$ multiple centres): stop climbing immediately; X-ray (may appear normal); MRI if X-ray negative and high suspicion; 6--12\,wk rest from climbing. Sch\"{o}ffl et al.\ (2025 BMJ Open) found physeal injuries in 29.9\% of 137 injuries across 95 adolescent patients. Applying adult pulley rehabilitation protocols to adolescents risks permanent growth-plate damage and joint deformity.
\end{warningbox}

%---------------------------------------------------------------------
\subsection{Skin Management}
\label{subsec:skin}

\subsubsection{Flappers --- Acute Traumatic Skin Tears}

A flapper is a traumatic partial-thickness skin tear, occurring at the callus-to-normal-skin transition where a stress riser under tangential shear load on rock hold edges causes delamination. Most common on the second and third digit palmar skin and thenar eminence.

\begin{folklorebox}{Acute Flapper Management: Universal Climbing Practice · \folkloreverified · ESS 1}
\textbf{Acute management (Folklore):} Trim any non-viable detached skin flap with sterile scissors (reduces bacterial trapping and maceration); irrigate; apply occlusive dressing or thin skin balm (petrolatum, Joshua Tree Salve, Climb On, or similar humectant) to maintain moist wound environment. Taping over a partially healed wound allows earlier return. \textbf{[Folklore: climbing community universal practice. General wound healing biology supports moist wound environment for epidermal regeneration (bioplausible extrapolation), but no climbing-specific trial exists.]}
\end{folklorebox}

\begin{toverifybox}{Return-to-Climbing Timeline After Flappers · \toverify · ESS 1}
\textbf{\toverify:} Anecdotal community consensus: mild flapper (superficial, $<$1\,cm$^2$) 3--7\,d; moderate (deeper or larger) 7--14\,d; larger raw dermis 1--2\,wk --- dependent on pain, infection risk, and hold texture. No empirical healing timeline data exist. Requires $\geq$3 independent named climbing sources with specific timelines for \folkloreverified status.
\end{toverifybox}

\begin{folklorebox}{Callus Management for Flapper Prevention · \folkloreverified · ESS 1}
\textbf{Prevention --- callus management (Folklore):} Filing calluses flush (reducing height and smoothing the callus-to-normal-skin transition) is near-universal in elite climbing practice; mechanistic goal is to reduce the stress riser at the thickness transition. Has not been evaluated in any study measuring flapper incidence as a function of callus management. Additional community recommendations: moisturise nightly, avoid alcohol-based sanitisers before climbing, manage session volume on sharp holds, rotate grip styles. \textbf{[Folklore: Eric Horst, \textit{Training for Climbing}; Lattice Training / Will Anglin; athlete and coach consensus.]}
\end{folklorebox}

\subsubsection{Chronic Skin Thinning}

\begin{bioplausbox}{Chronic Skin Thinning: Mechanism from Repeated Hold Friction · \bioplausible · ESS 2}
\textbf{Mechanism (Bioplausible):} Repeated friction on small-radius holds reduces stratum corneum thickness below functional baseline, presenting as hypersensitivity, split skin, and reduced friction coefficient. Contributing factors: training volume, hold substrate roughness, ambient humidity (very low humidity accelerates transepidermal water loss and skin cracking), and individual skin phenotype. No climbing-specific study has measured stratum corneum thickness longitudinally as a function of training load. \textbf{[Bioplausible: mechanism from general dermatology; no climbing-specific intervention data.]}
\end{bioplausbox}

\begin{folklorebox}{Climbing Skin Products and Humidity Management · \folkloreverified · ESS 1}
\textbf{Interventions (Folklore):} Climbing-specific skin products (Rhino Skin Solutions, ClimbSkin, Friction Labs products), humidity management (40--60\% indoor humidity anecdotally preferred), Antihydral (methenamine; acts by protein denaturation to reduce sweating, potentially increasing friction but risking ``glassy'' skin and splits), callus sanding, and session spacing are community-recommended. No peer-reviewed controlled trial has evaluated any climbing skin product for skin health, injury rate, or return-to-climbing outcomes. \textbf{[Folklore: manufacturer claims, community/practitioner consensus; Dave MacLeod, ClimbSkin. Emollient/humectant mechanisms are Bioplausible from general dermatology literature on skin barrier repair.]}
\end{folklorebox}

\subsubsection{Collagen Supplementation for Skin Healing}

\begin{evidencebox}{Collagen Supplementation: Skin and Connective Tissue Endpoints in Non-Climbing Populations · \evidencebacked · ESS 3}
\textbf{Collagen synthesis markers and skin endpoints (Evidence-backed for non-climbing populations):}
\begin{itemize}[leftmargin=1.5em]
  \item Shaw et al.\ (2017, \cite{SLBR+17}): RCT crossover ($n=8$); 15\,g gelatin + 48\,mg vitamin\,C, 1\,hr pre-exercise vs.\ placebo doubled circulating P1NP (collagen synthesis marker); in-vitro engineered ligament mechanics improved. \textit{Note: $n=8$ is very small; endpoints are surrogate; no injury or skin-healing outcome data.}
  \item A meta-analysis of oral hydrolyzed collagen peptides (26 RCTs, $n=1721$) for skin aging found improved hydration and elasticity endpoints. A systematic review/meta-analysis of 19 RCTs ($n=1341$) on oral/topical peptides found improved hydration and modest wrinkle reduction.
  \item Proksch et al.\ (2025, \cite{PZO25}): RCT ($n=66$), 2.5\,g collagen peptides vs.\ placebo --- reduced wrinkles, improved elasticity and hydration.
  \item A double-blind pilot RCT in burn patients ($n=30$): hydrolyzed collagen supplementation improved wound healing markers and reduced hospital stay $\approx$30\% vs.\ control.\singlesource
\end{itemize}
\textbf{[Evidence-backed for skin elasticity/hydration endpoints in aging and burn populations. Bioplausible for climbing skin healing: no direct evidence for flapper healing or friction-induced thinning in climbers.]}
\end{evidencebox}

%---------------------------------------------------------------------
\subsection{General Injury Prevention Principles}
\label{subsec:prevention}

\subsubsection{Warm-Up}

\begin{bioplausbox}{Warm-Up Efficacy for Injury Prevention: Mechanism Without Climbing RCT · \bioplausible · ESS 2}
\textbf{Warm-up efficacy (Bioplausible):} No climbing-specific RCT exists for injury prevention via warm-up. General sports medicine literature supports warm-up via increased connective tissue temperature (improving viscoelastic properties), increased neuromuscular activation, and enhanced motor pattern recruitment. The most methodologically robust proximate-domain evidence is the FIFA 11+ protocol (Soligard et al.\ 2008), which demonstrated $\approx$50\% lower-limb injury reduction in a large cluster-RCT in soccer --- tissue type, injury mechanism, and loading characteristics differ substantially from climbing finger pulley injuries.

Practical climbing warm-up: 5--10\,min aerobic movement; easy traversing; progressive finger recruitment; submaximal hangs before limit bouldering. \textbf{[Bioplausible: mechanistic basis exists; climbing-specific RCT evidence absent.]}
\end{bioplausbox}

\subsubsection{NSAIDs in Finger Injury Management}

\begin{bioplausbox}{NSAIDs for Finger Injuries: Contested Evidence, Pathology-Dependent Guidance · \bioplausible · ESS 2}
\textbf{NSAIDs (Bioplausible --- context-dependent):} NSAIDs are widely self-administered by climbers but their role in finger-injury healing is contested and pathology-dependent:
\begin{itemize}[leftmargin=1.5em,noitemsep]
  \item \textbf{Tenosynovitis / capsulitis (primarily inflammatory):} Short-course NSAIDs (48--72\,h acute phase) are consistent with anti-inflammatory management and broadly accepted in sports medicine. Evidence for clinical benefit in these conditions comes from general tendinopathy and synovitis literature, not climbing-specific trials.
  \item \textbf{Pulley / fibrocartilaginous tissue healing:} The evidence is conflicted. Heinemeier et al.\ (2017, \cite{HHM+17}) human tendinopathic tissue study found NSAIDs suppress COX-mediated prostaglandin E2 production that is required for tendon remodelling; the clinical significance at short doses is uncertain. Multiple in-vitro and animal models show COX-2 inhibition impairs tenocyte proliferation and collagen synthesis. However, no study has directly measured NSAID effects on A2 pulley healing in climbers.
  \item \textbf{Practical guidance:} Acetaminophen/paracetamol is preferred for analgesia during the first 2\,wk of acute pulley injury when bone-tendon healing is most active. If NSAIDs are used, limit to the initial 48--72\,h inflammatory phase only; avoid prolonged use during active tendon remodelling (weeks 1--6).
\end{itemize}
\textbf{[Bioplausible: mechanistic rationale from \cite{HHM+17}; no climbing-specific RCT on NSAID effects on pulley healing. Climbing community widely uses ibuprofen chronically through pain --- no outcome data exist.]}
\end{bioplausbox}

\subsubsection{Load Management --- The 10\% Rule}

\begin{toverifybox}{The 10\,\% Load Progression Rule Applied to Climbing · \toverify · ESS 1}
\textbf{\toverify --- 10\% rule (unvalidated in climbing):} The heuristic that weekly training load increases should not exceed 10\% of the prior week's load is widely cited in running medicine and climbing coaching. Its evidence base is weaker than commonly assumed: it derives primarily from retrospective epidemiological observation, not prospective RCT. Systematic reviews (e.g., Impellizzeri et al.\ 2020, \cite{IWC+20}) note that the specific 10\% threshold is arbitrary and inconsistently supported across sports. Gabbett (2016, \cite{Gab16}) and the acute:chronic workload ratio literature provide a more quantitative framework but remain primarily observational, contested in threshold values, and derived from team sports. For climbing, the threshold is additionally complicated because intensity, grip type, board angle, session density, and concurrent finger-training volume interact in ways not captured by a single linear loading metric. Requires $\geq$3 independent climbing-community sources citing this specific threshold for \folkloreverified status.
\end{toverifybox}

\subsubsection{Experience-Dependent Injury Risk}

\begin{evidencebox}{Experience-Dependent Finger Training Risk (Sjöman 2023) · \evidencebacked · ESS 5}
\textbf{Sjöman et al.\ (2023) --- Experience-dependent association (Evidence-backed):} Cross-sectional survey of $n=434$ Finnish climbers (\cite{SGJ23}) examined the association between route/bouldering-specific finger strength training (RFT) and sport/performance-impairing finger injury (SPIIF):
\begin{itemize}[leftmargin=1.5em,noitemsep]
  \item In climbers with $\geq$6\,yr experience: RFT \textbf{negatively} associated with SPIIF ($\chi^2=4.57$, $p=0.03$) --- structured training may be protective.
  \item In climbers with $<$6\,yr experience who had reached $\geq$7a (French sport): RFT \textbf{positively} associated with SPIIF ($\chi^2=4.61$, $p=0.03$) --- training in rapidly progressing novices may increase injury risk.
\end{itemize}
\textit{Quality constraints: cross-sectional design prohibits causal inference; self-reported injury and training data; Finnish climbing population; no dose-response data; no adjustment for confounders (training volume, hold type, warm-up, concurrent activities).}
\end{evidencebox}

\begin{bioplausbox}{Bidirectional Injury Risk Pattern: Mechanistic Interpretation · \bioplausible · ESS 2}
\textbf{Mechanistic interpretation (Bioplausible):} The bidirectional pattern is consistent with theoretical models of training-load injury risk in developing athletes (rapid capacity acquisition outpacing structural tissue adaptation) and the protective adaptation hypothesis in experienced athletes (accumulated structural adaptation of the fibrous flexor sheath, pulley hypertrophy, and neuromuscular co-contraction patterns). The mechanism for protection in experienced climbers has not been directly measured. \textbf{[Bioplausible for mechanism; Evidence-backed for the observed association in this dataset only.]}
\end{bioplausbox}

%---------------------------------------------------------------------
\subsection{What the Literature Cannot Yet Answer}
\label{subsec:injury-gaps}

\begin{enumerate}[leftmargin=2em]
  \item \textbf{Progressive loading vs.\ rest-only for Grade I--II A2 injuries.} Does structured hangboard rehabilitation outperform rest for time-to-return, re-injury rate, or ultrasound-confirmed healing in Grade I--II tears? No RCT exists; a moderate-sized trial ($n \geq 60$, Grade I--II, ultrasound-confirmed outcomes) is feasible.

  \item \textbf{Optimal immobilisation duration and type.} For Grade II--III A2 injuries, does 0--7\,d vs.\ 10--14\,d immobilisation (with standardised pulley protection) alter long-term TBD, return timeline, or re-injury rate? No controlled comparison exists.

  \item \textbf{H-taping clinical effectiveness.} Schweizer (2000, \cite{Sch00}) established small biomechanical reductions; Sch\"{o}ffl et al.\ (2007, \cite{SES+07}) reported acute crimp-strength improvement; Salas et al.\ (2022, \cite{SMT+23}) cadaveric study found no protective effect. Neither confirms reduced re-injury rate or faster healing in a controlled outcome study. Does H-taping change clinical outcomes, or does it act primarily via pain modulation and confidence? An RCT comparing PPS + H-tape versus PPS alone for clinical outcomes has not been conducted.

  \item \textbf{Ultrasound-guided return-to-climb criteria.} Does using serial TBD measurement to gate progression (e.g., TBD $<$2\,mm under resisted flexion before advancing grip position) reduce re-injury relative to symptom-only progression? The current TBD cutoff range of 1.9--5.1\,mm for complete rupture precludes standardised criteria.

  \item \textbf{Load-to-failure properties and healing kinetics of the human A2 pulley under progressive loading.} No in-vivo or ex-vivo study has quantified how incremental hangboard loading affects pulley tissue mechanics during healing. Tyler Nelson--style protocols lack any quantitative validation in the target tissue.

  \item \textbf{Collagen supplementation for pulley and skin healing.} Shaw et al.\ (2017) demonstrated elevated P1NP surrogate markers ($n=8$). No study has measured actual A2 pulley recovery time, re-injury rate, or imaging-based healing in climbers using gelatin + vitamin\,C. Similarly, no RCT has evaluated oral collagen peptides for flapper healing or chronic skin thinning under climbing-specific friction trauma.

  \item \textbf{Structural mediators of experience-dependent injury risk.} The Sjöman et al.\ (2023) bidirectional association is observational. A prospective cohort tracking novices from onset of structured finger training through 2--5\,yr, with serial ultrasound pulley morphology assessment, grip-force measurement, and load monitoring, would provide the first mechanistic evidence for whether pulley hypertrophy or neuromuscular adaptation mediates the protective effect observed in experienced climbers.

  \item \textbf{Pulley-protection splint vs.\ H-taping for clinical outcomes.} Schneeberger \& Schweizer (2016, \cite{SS16}) established PPS efficacy for TBD reduction and return-to-climbing, but no RCT has compared PPS + H-tape versus PPS alone versus H-tape alone for re-injury rate, return timeline, or patient-reported outcomes in Grade II--III injuries.

  \item \textbf{NSAID effects on pulley healing.} Climbers routinely self-medicate with ibuprofen through pulley pain. No study has measured the effect on healing outcomes in the A2 pulley (a fibrocartilaginous structure distinct from both tendon and bone). Whether short-course NSAIDs aid or harm healing in this specific tissue is unknown.

  \item \textbf{Tenosynovitis management: load reduction content vs.\ timing.} Mohn et al.\ (2022, \cite{MSM+22}) found no significant difference in outcomes based on specific therapy content; a prospective RCT comparing structured progressive loading versus unstructured load avoidance would clarify whether any specific rehabilitation element beyond ``reduce load'' adds clinical value.

  \item \textbf{Capsulitis natural history and OA progression.} The Sch\"{o}ffl (2025, \cite{SLL+25}) staged algorithm shows 72\% return to prior level, but long-term OA progression in those 28\% with decreased performance is not quantified. A 5--10\,yr cohort linking PIP synovitis episodes to imaging-confirmed OA development would establish whether aggressive early treatment is warranted.
\end{enumerate}
